% TODO proof-reading
\subsection{シンプルな計算関数}

シンプルな計算関数を続けましょう。

\begin{lstlisting}[style=customjava]
public class calc
{
	public static int half(int a)
	{
		return a/2;
	}
}
\end{lstlisting}

\INS{iconst\_2}命令が使用される場合の出力です：

\begin{lstlisting}
  public static int half(int);
    flags: ACC_PUBLIC, ACC_STATIC
    Code:
      stack=2, locals=1, args_size=1
         0: iload_0       
         1: iconst_2      
         2: idiv          
         3: ireturn       
\end{lstlisting}
         
\INS{iload\_0}は関数の0番目の引数を取り、それをスタックにプッシュします。

\INS{iconst\_2}は2をスタックにプッシュします。
これら2つの命令の実行後、スタックは次のようになります：

% FIXME: TikZ
\begin{lstlisting}
      +---+
TOS ->| 2 |
      +---+
      | a |
      +---+
\end{lstlisting}

\INS{idiv}は\ac{TOS}の2つの値を取り、一方を他方で割り、
結果を\ac{TOS}に残します：

% FIXME: TikZ
\begin{lstlisting}
      +--------+
TOS ->| result |
      +--------+
\end{lstlisting}

\INS{ireturn}はそれを取って返します。

倍精度浮動小数点数で続けましょう：

\begin{lstlisting}[style=customjava]
public class calc
{
	public static double half_double(double a)
	{
		return a/2.0;
	}
}
\end{lstlisting}

\begin{lstlisting}[caption=定数プール]
...
   #2 = Double             2.0d
...
\end{lstlisting}

\begin{lstlisting}
  public static double half_double(double);
    flags: ACC_PUBLIC, ACC_STATIC
    Code:
      stack=4, locals=2, args_size=1
         0: dload_0       
         1: ldc2_w        #2                  // double 2.0d
         4: ddiv          
         5: dreturn       
\end{lstlisting}

同じですが、\INS{ldc2\_w}命令が定数プールから定数2.0をロードするために使用されます。

また、他の3つの命令には\IT{d}プレフィックスがあり、
\IT{double}データ型の値で動作することを意味します。

今度は2つの引数を持つ関数を使ってみましょう：

\begin{lstlisting}[style=customjava]
public class calc
{
	public static int sum(int a, int b)
	{
		return a+b;
	}
}
\end{lstlisting}

\begin{lstlisting}
  public static int sum(int, int);
    flags: ACC_PUBLIC, ACC_STATIC
    Code:
      stack=2, locals=2, args_size=2
         0: iload_0       
         1: iload_1       
         2: iadd          
         3: ireturn       
\end{lstlisting}

\INS{iload\_0}は最初の関数引数（a）をロードし、\INS{iload\_1}は2番目（b）をロードします。

両方の命令の実行後のスタックの状態です：

\begin{lstlisting}
      +---+
TOS ->| b |
      +---+
      | a |
      +---+
\end{lstlisting}

\INS{iadd}は2つの値を加算し、結果を\ac{TOS}に残します：

\begin{lstlisting}
      +--------+
TOS ->| result |
      +--------+
\end{lstlisting}

この例を\IT{long}データ型に拡張してみましょう：

\begin{lstlisting}[style=customjava]
	public static long lsum(long a, long b)
	{
		return a+b;
	}
\end{lstlisting}

\dots 次のような結果が得られます：

\begin{lstlisting}
  public static long lsum(long, long);
    flags: ACC_PUBLIC, ACC_STATIC
    Code:
      stack=4, locals=4, args_size=2
         0: lload_0       
         1: lload_2       
         2: ladd          
         3: lreturn       
\end{lstlisting}

2番目の\TT{lload}命令は2番目の引数を2番目のスロットから取得します。

これは、64ビット\IT{long}値が正確に2つの32ビットスロットを占有するためです。

少し高度な例：

\begin{lstlisting}[style=customjava]
public class calc
{
	public static int mult_add(int a, int b, int c)
	{
		return a*b+c;
	}
}
\end{lstlisting}

\begin{lstlisting}
  public static int mult_add(int, int, int);
    flags: ACC_PUBLIC, ACC_STATIC
    Code:
      stack=2, locals=3, args_size=3
         0: iload_0       
         1: iload_1       
         2: imul          
         3: iload_2       
         4: iadd          
         5: ireturn       
\end{lstlisting}

最初のステップは乗算です。積は\ac{TOS}に残されます：

\begin{lstlisting}
      +---------+
TOS ->| product |
      +---------+
\end{lstlisting}

\TT{iload\_2}は3番目の引数（c）をスタックにロードします：

\begin{lstlisting}
      +---------+
TOS ->|    c    |
      +---------+
      | product |
      +---------+
\end{lstlisting}

今度は\TT{iadd}命令が2つの値を加算できます。